\documentclass[11pt]{article}
\usepackage{amsmath,url}
\title{How Much Mule Preparation Can Wait During a Partition?}
\date{25 November 2026}
\begin{document}
\maketitle
\paragraph{Result.}
A rover retains an inspection report when its path to the base station
fails. Its first recovery assigns delivery to a relay; this recovery is
not settled until a named durable receipt gets back. A second fault
during that wait can remove the relay's route and transfer the duty to
a data mule. There is an exact minimum amount of mule staging to do
while no resolving observation has arrived. In the development
example, waiting for notice misses the delivery deadline under either
admitted second fault. Eager staging works but spends two preparatory
ticks if the first relay succeeds at tick 2. The minimum policy
spends one tick on that branch and still meets both second-fault deadlines.

\paragraph{Source and model.}
NeBula describes autonomous operation beyond its mesh, communication
with a returning data mule or backtracking to the mesh, and identification
of a restored heartbeat apart from motion and mission planning
(\url{https://arxiv.org/html/2103.11470}, sections 3, 10.1, and 11.2).
The model uses those operating options; it is not a run of NeBula,
CHORD, OPTIC, or a measured network. It adds a durable receipt
for a named report and keeps the producer's copy until that receipt.
The original relay can lose the ability to finish yet retain a live
motion authority at a shared copy station. Neither a lost heartbeat
nor a reported failure supplies evidence that the relay cannot enter.
The plant therefore supplies separate, durable station clearance,
old-entry exclusion, and exclusive access for the mule. The
observation that changes the supplier is different from each of these
three proofs. The no-entry premise in the openTCS publication replay
is one example of a boundary the model must respect, not an assumption
that those systems are deployed together.

Staging consumes one reserve work slot per tick along a private
approach. Partial staging persists. No staging action observes
the hidden failure before the resolving notice or alters when a
notice arrives. There is no monetary or physical field estimate;
the counted charge is staging work before the first resolving event,
not idle battery consumption. The receiving station is not entered
before the new grant. Independent inspection ticks are on a disjoint
footprint, without waiting on a global planner. Beyond the stated
fault family, late or absent evidence leaves the original report
kept and no new shared-station operation admitted. The prototype
has no queueing or deadline overlap with a second report, no common
loss of the retained copy, and no accidental reentry after the gate.
Those would change the declared envelope.

\paragraph{Exact preparation bound.}
Fix total staging work $w$ and a deadline $D$. In a failure branch
$h$, local notice arrives at $t_h$, separate gate/clearance/grant
facts arrive at $g_h,q_h,a_h$, and $d_h>0$ bounds the protected
copy-and-return operation, including delivery and egress.
These bounds apply to a single common history. The conservative
station claim lasts for that whole operation; the model does not
release it midway. If $x$ staging ticks had already been used
when notice arrived, the earliest certified finish is
\[
 \max(t_h+w-x,\ g_h,\ q_h,\ a_h)+d_h.
\]
No amount of preparation can compensate if a required fact arrives
after $D-d_h$. For other branches set
$r_h=\max(0,t_h+w+d_h-D)$, the minimum staging needed by $t_h$.
A feasible silent policy is nondecreasing, advances by at most one
per tick and never exceeds $w$. Thus if $r_h>\min(w,t_h)$
in any admitted case, the service floor is impossible within the model.

Otherwise the unique pointwise minimum cumulative staging on the
no-notice path is
\[
 P(u)=\max\bigl(0,\ \max_h[r_h-\max(t_h-u,0)]\bigr).
\]
For a future cut $t_h$, a policy can add at most $t_h-u$ ticks
after $u$; for a past cut its prior staging is nonrefundable. These
give the pointwise lower bound. Each term is nondecreasing and
grows by at most one tick per tick, as does their maximum.
The maximum starts at zero, stays at most $w$, and provides the
required staging at each notice. Perform a staging tick exactly
when the maximum increases and stop speculative staging when a
resolving observation arrives. In a failure branch, resume just
the remaining work, and use the frozen receiver with the old entry,
clearance and grant guarantees. This proves both safety and
the service floor under the stated physical premises and every
branch in the declared family.

\paragraph{Matched checks.}
The calibration has $w=4$, $D=8$, notice times $2,4$, and protected
duration 3. The two last required fact times are 4 and 5. Minimum
staging by ticks 2 and 4 is one and three; finish is 8 in both
failure branches. Moving the second branch's final required fact from tick 5 to
tick 6 breaks the deadline even with full staging. A
direct finite policy search with the same actions and observations
matches pointwise, including infeasibility, in 1,500 enumerated
two-outcome cases. The proof's value here is characterizing what
can wait, not outperforming such a policy. Eager staging and delayed
replanning are ablations, not substitutes for the strongest comparator.
Lean checks the integer staging requirement and the lemmas for the
pointwise bound; the Python decision code and its source adapter are
not verified in Lean. This is one development operating family,
not an independently selected mission or an obligation acceptance.
\end{document}
